\documentclass[11pt,a4paper]{report}
% Préambule partagé des rapports SecuriGlass Factory OS.
% Style volontairement sobre : pas d'encadrés de couleur, pas de blocs
% décoratifs. Règles horizontales fines, typographie classique, une seule
% couleur d'accent reprise de l'identité visuelle du produit.

\usepackage{fontspec}
\usepackage{polyglossia}
\setdefaultlanguage{french}
\setmainfont{Palatino}
\setsansfont{Helvetica Neue}
\setmonofont{Menlo}[Scale=0.85]
\usepackage{microtype}
\sloppy
\usepackage[a4paper,top=2.8cm,bottom=2.8cm,left=2.6cm,right=2.6cm,headsep=1.1cm]{geometry}
\usepackage{graphicx}
\usepackage{booktabs}
\usepackage{longtable}
\usepackage{array}
\usepackage{caption}
\usepackage{enumitem}
\usepackage{xcolor}
\usepackage{fancyhdr}
\usepackage{titlesec}
\usepackage{listings}
\usepackage{tocloft}
\usepackage[hidelinks]{hyperref}
\usepackage[nameinlink,french]{cleveref}

% ---------------------------------------------------------------- couleurs
\definecolor{accent}{HTML}{15584B}
\definecolor{accentlight}{HTML}{1F7A68}
\definecolor{graytext}{HTML}{3F4A47}
\definecolor{rulegray}{HTML}{C9CFCD}

\hypersetup{
  colorlinks=true,
  linkcolor=accent,
  citecolor=accent,
  urlcolor=accentlight,
  pdfauthor={Mohamed Aziz Baoueb},
}

% ------------------------------------------------------------------ titres
\titleformat{\chapter}[hang]
  {\normalfont\Huge\bfseries\color{accent}}
  {\thechapter}{1.2cm}{}
\titlespacing*{\chapter}{0pt}{-0.8cm}{1.0cm}

\titleformat{\section}
  {\normalfont\Large\bfseries\color{accent}}
  {\thesection}{0.6em}{}
\titleformat{\subsection}
  {\normalfont\large\bfseries\color{graytext}}
  {\thesubsection}{0.6em}{}
\titleformat{\subsubsection}
  {\normalfont\normalsize\bfseries\itshape\color{graytext}}
  {\thesubsubsection}{0.6em}{}

% --------------------------------------------------------------- en-têtes
\pagestyle{fancy}
\fancyhf{}
\renewcommand{\headrulewidth}{0.4pt}
\renewcommand{\footrulewidth}{0pt}
\fancyhead[L]{\small\itshape\nouppercase{\leftmark}}
\fancyhead[R]{\small\thepage}
\fancypagestyle{plain}{
  \fancyhf{}
  \fancyhead[R]{\small\thepage}
  \renewcommand{\headrulewidth}{0.4pt}
}

% ----------------------------------------------------------------- listes
\setlist{itemsep=2pt, topsep=4pt, parsep=0pt}

% --------------------------------------------------------------- légendes
\captionsetup{
  labelfont={bf,color=accent},
  font=small,
  skip=6pt,
  justification=centering,
}

% ----------------------------------------------------------- code source
\lstset{
  basicstyle=\ttfamily\small,
  breaklines=true,
  frame=none,
  columns=flexible,
  keywordstyle=\color{accent}\bfseries,
  commentstyle=\color{graytext}\itshape,
  stringstyle=\color{accentlight},
  showstringspaces=false,
  tabsize=2,
}

% -------------------------------------------------------- table des matières
\renewcommand{\cftchapfont}{\bfseries\color{accent}}
\renewcommand{\cftchappagefont}{\color{accent}}

% ----------------------------------------------------------- environnement
% figure pleine largeur avec légende, sans cadre ni fond.
\newcommand{\screenshot}[3]{%
  \begin{figure}[htbp]
    \centering
    \includegraphics[width=#2\textwidth]{figures/#1}
    \caption{#3}
  \end{figure}%
}

% séparateur fin entre deux sections d'un même thème (jamais de remplissage).
\newcommand{\finesep}{%
  \par\noindent\textcolor{rulegray}{\rule{\linewidth}{0.4pt}}\par\vspace{4pt}%
}


\title{SecuriGlass Factory OS}
\author{Mohamed Aziz Baoueb}
\date{Octobre 2026}

\begin{document}

% ============================================================ page de titre
\begin{titlepage}
  \centering
  \vspace*{2.2cm}

  {\Large\scshape Rapport de projet}\\[0.6cm]
  \textcolor{rulegray}{\rule{0.5\linewidth}{0.4pt}}\\[1.4cm]

  {\fontsize{32}{38}\selectfont\bfseries\color{accent} SecuriGlass Factory~OS}\\[0.5cm]
  {\Large\itshape Console de supervision sécurisée pour une ligne\\
  de production de verre trempé}\\[2.2cm]

  \begin{tabular}{rl}
    \textbf{Auteur} & Mohamed Aziz Baoueb \\[4pt]
    \textbf{Date} & Octobre 2026 \\[4pt]
    \textbf{Nature} & Preuve de concept (PoC) \\[4pt]
    \textbf{Domaine} & Cybersécurité industrielle (OT/IT), IEC~62443, OWASP ASVS \\
  \end{tabular}

  \vfill
  \textcolor{rulegray}{\rule{0.5\linewidth}{0.4pt}}\\[0.4cm]
  {\small Référentiels appliqués : IEC~62443-3-3, OWASP ASVS v4.0, OWASP Top~10 for LLM,\\
  NIST SP~800-82, RGPD, norme EN~12150-1 (verre de sécurité).}
\end{titlepage}

\tableofcontents

% ================================================================ chapitre
\chapter{Contexte et objectifs}

\section{Pourquoi ce projet}
SecuriGlass Factory~OS est une preuve de concept (PoC) qui simule la supervision
numérique d'une ligne de production de verre trempé, depuis l'automate qui pilote
le four jusqu'au tableau de bord consulté par les équipes de l'usine. L'objectif
n'est pas de livrer un produit fini, mais de démontrer, sur un cas industriel
concret, comment concevoir un système qui relie un réseau industriel (OT) à un
réseau informatique (IT) \emph{sans} affaiblir la sécurité de l'un ou de l'autre.

Le verre trempé est un bon cas d'étude : sa fabrication suit une norme physique
stricte (EN~12150-1, qui fixe les plages de température et de pression de
trempe), et une dérive du four n'est pas qu'un problème informatique — elle peut
casser la production ou, dans un cas réel, mettre en danger une personne. Le
projet traite donc la sécurité informatique et la sécurité physique du procédé
comme une seule et même préoccupation.

\section{Objectifs poursuivis}
\begin{itemize}
  \item Construire une chaîne de bout en bout réaliste : automate simulé en C,
    passerelle de collecte, liaison unidirectionnelle vers le système
    d'information, base de données, API, console web.
  \item Appliquer le modèle de Purdue (zonage par niveaux) et le principe de
    moindre privilège à chaque frontière du système.
  \item Donner à un assistant conversationnel (IA) un accès utile aux données de
    production, sans lui laisser la possibilité de divulguer des informations
    sensibles ni d'agir sur l'automate.
  \item Démontrer, par des tests d'intrusion réels et non scénarisés, que les
    protections tiennent — et documenter honnêtement ce qui ne tient pas encore.
  \item Produire une console que des rôles différents (direction d'usine,
    opérateur, qualité, audit sécurité) peuvent utiliser sans qu'aucun ne voie
    plus que ce dont il a besoin.
\end{itemize}

\section{Périmètre et limites assumées}
Ce projet est une démonstration technique, pas un produit industriel homologué.
Certaines mesures citées dans la documentation (ex. réplication à froid,
signature des consignes à l'automate, image Docker \emph{distroless}) sont
volontairement marquées comme \emph{planifiées} plutôt qu'implémentées, pour ne
jamais laisser croire qu'une protection existe alors qu'elle n'a pas été codée.
Cette honnêteté sur le périmètre fait partie de la méthode : un rapport de
sécurité qui surestime ses garanties est plus dangereux qu'un rapport qui les
sous-estime.

\chapter{Vue d'ensemble du produit}

\section{Les six rôles de la console}
L'accès à la console est entièrement déterminé par le rôle du compte connecté,
aussi bien dans l'interface que côté serveur (une tentative d'accès à une
ressource non autorisée échoue même si l'utilisateur en devine l'adresse).

\begin{center}
\begin{tabular}{@{}ll@{}}
\toprule
\textbf{Rôle} & \textbf{Ce qu'il voit} \\
\midrule
Super administrateur & L'ensemble de la console \\
Responsable d'usine & Production, stocks, commandes, clients, assistant \\
Opérateur de ligne & Production, stocks, assistant \\
Inspecteur qualité & Production uniquement \\
Auditeur sécurité & Sécurité, journal d'audit, clients \\
Compte de service (passerelle) & Aucun accès à la console \\
\bottomrule
\end{tabular}
\end{center}

\section{Parcours d'une commande, de la vente à la trempe}
Pour donner une vue concrète du système, voici ce qui se passe quand un client
commande des vitres :
\begin{enumerate}
  \item Le responsable d'usine enregistre la commande dans la console ; le
    serveur vérifie qu'il reste assez de matière première, puis déduit
    automatiquement les quantités nécessaires du stock.
  \item Le four, piloté par l'automate (simulé en C), chauffe le verre selon la
    recette choisie ; les mesures (température par zone, pression de trempe,
    vitesse de convoyeur) sont relevées en continu.
  \item Un modèle d'apprentissage automatique (forêt d'isolement) surveille ces
    mesures et signale tout écart par rapport au comportement normal du four.
  \item Les mesures et les éventuelles alertes sont visibles en direct dans la
    console, avec un diagnostic écrit en français.
  \item Un opérateur peut interroger l'assistant technique pour obtenir le
    rapport de production ou une explication de la norme EN~12150 ; l'assistant
    n'a accès à aucune donnée financière ou personnelle des clients.
  \item Toute action sensible (connexion, création d'un client, tentative
    d'accès refusée, consigne bloquée) est inscrite dans un journal d'audit
    scellé, consultable par l'auditeur sécurité.
\end{enumerate}

\chapter{Fonctionnalités}

\section{Accueil}
La page d'accueil résume en un coup d'œil ce qui concerne le rôle connecté :
température du four, état du four selon l'analyse automatique, lot en cours,
stocks à commander, commandes en cours, et pour les rôles concernés la note de
conformité et le nombre d'attaques bloquées sur les dernières 24~heures.

\screenshot{02_accueil.jpg}{0.92}{Page d'accueil, vue du responsable d'usine.}

\section{Production}
Cette page affiche la température du four en direct, zone par zone, avec une
courbe qui change de couleur selon le diagnostic du modèle d'apprentissage
automatique : bleu pour un fonctionnement normal, orange pour un écart à
surveiller, rouge pour une anomalie. Les jauges de température et de pression de
soufflage complètent la lecture.

\screenshot{03_production.jpg}{0.92}{Suivi du four pendant un cycle de chauffe réel ; le modèle signale un écart entre zones.}

\section{Stocks}
Les six matières premières nécessaires à la fabrication (sable, carbonate de
sodium, calcaire, calcin recyclé, revêtement bas-émissif, palettes) sont
affichées avec leur niveau. Un stock passé sous son seuil de réapprovisionnement
est signalé automatiquement.

\screenshot{04_stocks.jpg}{0.92}{Page des stocks de matières premières.}

\section{Commandes}
Le responsable d'usine enregistre une commande en choisissant un client, une
recette de verre et une quantité. Le serveur calcule la surface de verre, vérifie
qu'il y a assez de matière, l'enregistre, puis déduit automatiquement les
quantités consommées des stocks — le tout dans une seule opération qui annule
tout en cas de problème.

\screenshot{05_commandes.jpg}{0.92}{Prise de commande et historique des commandes récentes.}

\section{Clients}
Le fichier clients contient des données personnelles (nom du contact, e-mail,
téléphone) chiffrées avant d'être écrites en base ; seule la raison sociale
reste en clair pour permettre la recherche. L'accès en création est réservé au
responsable d'usine.

\screenshot{06_clients.jpg}{0.92}{Fichier clients et formulaire de création.}

\section{Assistant technique}
L'assistant répond aux questions sur la production et sur la norme EN~12150. Il
est protégé par quatre barrières successives (filtrage de la question, réponses
puisées dans une base documentaire fermée, génération contrainte, vérification
des grandeurs physiques en sortie) qui l'empêchent de divulguer des informations
internes ou d'accepter une instruction cachée dans la question de l'utilisateur.
La capture ci-dessous montre une tentative réelle d'injection de prompt,
bloquée par le filtre et consignée dans le journal de sécurité.

\screenshot{07_assistant.jpg}{0.92}{L'assistant refuse une tentative d'injection de prompt en direct.}

\section{Sécurité}
Réservée aux rôles habilités, cette page regroupe la note de conformité (part
des contrôles de sécurité qui sont en règle), le flux des tentatives bloquées
qualifiées selon le référentiel international MITRE~ATT\&CK, et le schéma de
circulation des données entre le réseau industriel et le réseau informatique.

\screenshot{08_securite.jpg}{0.92}{Tableau de conformité et premières alertes de sécurité.}
\screenshot{09_securite_reseau.jpg}{0.92}{Flux d'alertes qualifiées MITRE et schéma de circulation des données OT~$\rightarrow$~IT.}

\section{Journal d'audit}
Chaque action sensible est inscrite dans un journal chaîné par empreintes
SHA-256 : chaque entrée contient l'empreinte de la précédente, ce qui rend toute
modification rétroactive détectable. Un bouton recalcule la chaîne entière côté
serveur pour en vérifier l'intégrité.

\screenshot{10_audit.jpg}{0.92}{Journal d'audit et résultat de la vérification de la chaîne.}

\chapter{Architecture et sécurité}

\section{Le modèle de Purdue, du four au navigateur}
Le système est organisé en zones de confiance croissante, conformément au
modèle de Purdue utilisé en cybersécurité industrielle :

\begin{center}
\begin{tabular}{@{}cll@{}}
\toprule
\textbf{Niveau} & \textbf{Zone} & \textbf{Composant} \\
\midrule
0--1 & Terrain / automate & Automate du four (simulateur C, Modbus~TCP) \\
2--3 & Supervision OT & Passerelle de collecte (lecture seule) \\
3,5 & Zone intermédiaire & Diode de données (liaison à sens unique) \\
4 & Informatique de gestion & API, base de données \\
5 & Entreprise & Console web, assistant, journal d'audit \\
\bottomrule
\end{tabular}
\end{center}

Le point central de cette architecture est la \textbf{diode de données} : la
passerelle ne fait qu'émettre des trames UDP vers le système d'information, sans
jamais recevoir de réponse. Il n'existe, techniquement, aucun chemin réseau pour
qu'une commande venant du portail web atteigne l'automate. C'est ce qui a été
vérifié en pratique (\cref{chap:resultats}).

\section{Défense en profondeur}
\begin{itemize}
  \item \textbf{Identité} : mots de passe en Argon2id, jetons de session signés
    et révocables, verrouillage après plusieurs échecs.
  \item \textbf{Autorisation} : contrôle par rôle (RBAC) sur chaque route de
    l'API, complété par des règles contextuelles (ABAC) refusant tout par
    défaut.
  \item \textbf{Confidentialité} : les données personnelles des clients sont
    chiffrées au repos (AES-256-GCM) ; en cas de fuite de la base, elles restent
    illisibles.
  \item \textbf{Intégrité des mesures} : chaque trame envoyée par la passerelle
    est scellée (empreinte SHA-256 et signature Ed25519) et numérotée pour
    empêcher le rejeu.
  \item \textbf{Assistant IA} : quatre barrières successives, détaillées au
    \cref{chap:technique} du rapport technique.
  \item \textbf{Traçabilité} : journal d'audit chaîné et immuable, y compris
    pour les tentatives refusées.
\end{itemize}

\section{Ce que l'architecture ne couvre pas encore}
Une attaque réseau directe sur l'automate (en contournant le portail web) n'est
aujourd'hui pas détectée automatiquement : le flux d'alertes ne surveille que ce
qui passe par l'application, pas le réseau industriel lui-même. Cette limite a
été découverte pendant les tests d'intrusion (\cref{chap:resultats}) et n'a pas
été dissimulée.

\chapter{Démarche de développement}

\section{Une progression par étapes}
Le projet a été construit en sept étapes successives, chacune validée avant de
passer à la suivante : fondations données et identité, simulation du réseau
industriel, services métier et intelligence artificielle, interface web,
supervision de sécurité, infrastructure, puis homologation documentaire. Cette
progression a permis de ne jamais construire une couche sur des fondations non
vérifiées.

\section{Qualité et vérification continue}
Chaque changement est passé par les mêmes portes de qualité, automatisées dans
une intégration continue : vérification des types, analyse statique de
sécurité, formatage, tests, audit des dépendances, construction des images
Docker, génération d'un inventaire logiciel (SBOM) et signature des artefacts.
Le détail chiffré figure au \cref{chap:resultats}.

\section{Itérations sur l'interface}
L'interface a connu plusieurs itérations de style avant sa version actuelle,
chacune guidée par un retour direct sur l'usage : d'abord une console sombre
très contrastée, jugée trop proche d'un jeu vidéo ; puis une version claire et
sobre, jugée trop froide ; la version retenue reprend un principe de panneaux
translucides sur fond clair, avec des transitions discrètes, un mode sombre qui
suit la préférence du système, et une bascule complète français/anglais.

\chapter{Résultats et vérification}
\label{chap:resultats}

\section{Ce qui a été mesuré}
\begin{center}
\begin{tabular}{@{}ll@{}}
\toprule
\textbf{Composant} & \textbf{Résultat} \\
\midrule
API et logique métier (Python) & 563 tests passés \\
Interface web (TypeScript/React) & 122 tests passés \\
Automate (C) & 22 tests unitaires, 763 cas de fuzzing, 0 fuite mémoire \\
Passerelle de collecte (Python) & 20 tests passés \\
Typage statique strict & 0 erreur \\
Analyse de sécurité statique & 0 anomalie de sévérité moyenne ou haute \\
Intégration continue & 7 étapes automatisées (test, sécurité, image, signature) \\
\bottomrule
\end{tabular}
\end{center}

\section{Une attaque réelle, pas un scénario écrit à l'avance}
Pour vérifier que les protections tenaient vraiment, une attaque réseau a été
menée directement contre l'automate simulé, sans passer par le portail web, à
l'aide d'un outil de test d'intrusion dédié au projet :
\begin{enumerate}
  \item Une consigne de température absurde (900\textdegree{}C) a été
    \textbf{refusée par l'automate lui-même}, avec un code d'erreur réseau
    normal — ce n'est pas un texte préécrit, c'est une limite programmée dans
    son micrologiciel.
  \item Une consigne dans les bornes légales mais dangereuse (750\textdegree{}C
    sur deux zones sur quatre) a été \textbf{acceptée}. Le four a réellement
    chauffé au-delà de la plage de sécurité affichée (620--730\textdegree{}C),
    et le modèle d'apprentissage automatique a détecté l'écart \textbf{sans
    aucune aide}, en continu.
  \item Un arrêt d'urgence envoyé par le réseau a coupé le chauffage, mais il
    s'est avéré \textbf{impossible de le lever à distance} : comme sur une
    machine réelle, seule une intervention physique le permet.
\end{enumerate}
Cette expérience a aussi révélé la limite citée au chapitre précédent : cette
attaque-là, menée directement sur le réseau industriel, n'apparaît pas dans le
flux d'alertes de la console, faute de sonde sur ce réseau.

\chapter{Difficultés rencontrées}

\section{Collisions entre styles et objectifs}
Le principal défi n'a pas été technique mais une tension d'objectifs : rendre
l'interface à la fois sobre (pour un utilisateur professionnel de cinquante
ans), vivante (pour ne pas paraître austère) et parfaitement lisible sans
jamais verser dans l'esthétique « jeu vidéo ». Plusieurs itérations complètes de
la feuille de style ont été nécessaires avant de trouver un équilibre jugé
satisfaisant.

\section{Du CSS hérité qui reprenait la main}
Un bogue d'affichage (menu qui redevenait horizontal, barre du haut qui
débordait) provenait de règles CSS héritées d'une version antérieure du thème,
toujours actives sur certaines largeurs d'écran bien qu'invisibles au premier
regard. La cause a été trouvée en comparant le comportement calculé par le
navigateur à la cascade de règles réellement appliquées, plutôt qu'en devinant.

\section{Un modèle trop prudent}
Le détecteur d'anomalies a d'abord signalé presque toutes les mesures comme
suspectes : il n'avait jamais vu de montée en température ni de four à l'arrêt
dans ses données d'entraînement. Le jeu d'entraînement a été enrichi pour
couvrir ces régimes de fonctionnement normaux, ramenant le taux de fausses
alertes à un niveau acceptable sans perdre la détection des pannes réelles.

\chapter{Conclusion et perspectives}

SecuriGlass Factory~OS démontre, sur un cas industriel concret, qu'il est
possible de relier un procédé physique critique à une console web moderne sans
sacrifier la séparation stricte entre réseau industriel et réseau
informatique — et que cette séparation se vérifie par l'expérience, pas
seulement sur le papier. Le projet reste une preuve de concept~: certaines
protections citées dans la documentation technique sont volontairement
marquées comme planifiées, et l'homologation n'a pas été prononcée.

Les suites naturelles du projet sont, par ordre de priorité~: ajouter une
détection des attaques réseau directes sur l'automate (constatée manquante
pendant les tests), faire tourner la chaîne de bout en bout sur un système
Linux pour exécuter les tests aujourd'hui ignorés faute d'outillage macOS, puis
traiter les mesures marquées « planifiées » dans le dossier d'architecture.

\end{document}
